\newpage
\phantomsection
\label{prf:equiv-partition}
\LRAProofFor{thm:equiv-partition}

\begin{remark*}[Return]
\hyperref[thm:equiv-partition]{Return to Theorem}
\end{remark*}

\begin{theorem*}[Equivalence Relations and Partitions]
Let $A$ be a set.
\begin{enumerate}[label=(\roman*)]
\item If $R$ is an equivalence relation on $A$, then $A/R$ is a partition of $A$.
\item If $\mathcal{P}$ is a partition of $A$, then the relation
$R_{\mathcal{P}}$ defined by
\[
(a,b) \in R_{\mathcal{P}} \;\Longleftrightarrow\;
\exists P \in \mathcal{P} \text{ with } a \in P \text{ and } b \in P
\]
is an equivalence relation on $A$.
\item These constructions are inverse: $R_{A/R} = R$ and $A/R_{\mathcal{P}} = \mathcal{P}$.
\end{enumerate}
\end{theorem*}

\begin{proof}[Professional Standard Proof]
\LRAProofBodyStart
TODO: professional standard proof for equiv-partition.
\end{proof}

\begin{proof}[Detailed Learning Proof]
\LRAProofBodyStart
TODO: detailed learning proof for equiv-partition.
\end{proof}

\begin{remark*}[Proof structure]
\end{remark*}

\begin{dependencies}
\begin{itemize}
  \item TODO: list mathematical dependencies.
\end{itemize}
\end{dependencies}

\clearpage
